\documentclass[10pt]{article}
% -------------------------------------------------------------------
\usepackage[english]{babel}										
\usepackage[utf8]{inputenc}										% Allows for accented characters
\usepackage[T1]{fontenc}										% Accented and other characters
\usepackage{hyperref}
% -------------------------------------------------------------------
% Math packages
\usepackage{amsmath,amsfonts,amssymb,amsthm,cancel,siunitx,
calculator,calc,mathtools,empheq,latexsym}
% -------------------------------------------------------------------
% Packages for figures
\usepackage{subfig,epsfig,tikz,float}		            % For Figures
% -------------------------------------------------------------------
% Packages for tables and bibliography
\usepackage{booktabs,multicol,multirow,tabularx,array}          % For tables
\usepackage{natbib}
% -------------------------------------------------------------------
% Define the layout of the document
\setlength{\parindent}{0pt}
\setlength{\parskip}{5pt}
\textwidth 13.5cm
\textheight 19.5cm
\columnsep .5cm

% -------------------------------------------------------------------
% Article title
% \title{\renewcommand{\baselinestretch}{1.17}\normalsize\bf%

\title{\renewcommand{\baselinestretch}{1.17}\bf
{Inferring the Political Weights of Presidential Candidates from Their Tax Policy Proposals}
}
% -------------------------------------------------------------------
% Authors
\author{
Jason DeBacker\thanks{Jason DeBacker, Darla Moore School of Business, University of South Carolina,  \href{mailto:jason.debacker@moore.sc.edu}{jason.debacker@moore.sc.edu}} \ and \ John P. Ryan\thanks{John P. Ryan, University of Wisconsin,  \href{mailto:john.p.ryan@wisc.edu}{john.p.ryan@wisc.edu}}
}
% -------------------------------------------------------------------

%The document

\begin{document}

\date{}

\maketitle

\vspace{-0.5cm}


% -------------------------------------------------------------------
% Abstract
\bigskip
\noindent
{\small{\bf ABSTRACT.}
We use the method of inverse optimal taxation to infer political weights from the tax policy proposals of recent presidential candidates in the United States.  This method allows us to recover the implicit weights these candidates place on taxpayers across the income distribution. We find generally decreasing weights beyond modal income with significant differences between candidates and parties at the high end of the income distribution. Democratic candidates place lower weight on high-income taxpayers than Republican candidates. From 2012 to 2024, we note that both parties reduce the political weight they attribute to taxpayers with incomes above \$400,000. We also ask whether differences in policy proposals can be explained by heterogeneous beliefs about the behavioral responses to taxes, and find that modest differences in beliefs about taxpayer responses can explain significant differences in policies.
}

\medskip
\noindent
{\small{\bf Keywords}{:} 
Inverse Optimal Taxation, Optimal Taxation, Political Preferences, Tax Policy
}

\baselineskip=\normalbaselineskip
% -------------------------------------------------------------------

\section{Introduction}\label{sec:1}

What do presidential candidates' tax proposals reveal about their underlying social preferences? This paper answers that question using the inverse-optimal tax framework, recovering the implicit political weights that U.S. presidential candidates assign to taxpayers across the income distribution. Tax policy is among the most contested terrain in presidential campaigns — candidates debate the appropriate degree of progressivity, the revenue needed to finance other priorities, and the incentive effects of taxation — but these debates rarely make explicit the distributional value judgments that motivate each side. Inverse optimal taxation, a revealed preference methodology, makes those implicit judgments explicit.

We apply the approach developed by \citet{JJZ2017} — who estimate social welfare weights from Dutch tax policy — to U.S. presidential candidates' tax proposals since 2012, combining those proposals with empirical estimates of the income distribution and the elasticity of taxable income (ETI) to recover each candidate's implicit political weights. Our application extends this framework to a multi-candidate, partisan comparison setting, allowing us to characterize how distributional preferences differ across parties and election cycles.

We find that Democratic and Republican candidates alike exhibit a high degree of status-quo bias — their proposed tax schedules imply political weights that differ only modestly from those embedded in current law across most of the income distribution. Across all candidates, political weights decline with income, reflecting a shared preference for progressivity. The sharpest partisan divide appears at the top: Democratic candidates consistently assign lower weight to very high-income taxpayers than their Republican counterparts, though candidates from both parties assign similar weights to those below the 90th percentile. Interestingly, Trump and Clinton's 2016 proposals implied nearly identical weights for taxpayers earning up to \$650,000, suggesting some compression of partisan differences relative to 2012.

Beyond documenting these patterns, we examine how assumptions about behavioral responses shape the inferred weights. Under a constant ETI, political weights at the top of the distribution are highly sensitive to that assumption — higher elasticities imply lower weights on top earners under the progressive schedules all candidates proposed. When we instead interpolate income-varying ETIs from the empirical literature, the implied weights on high-income taxpayers fall substantially relative to the constant-elasticity baseline. Most strikingly, we show that differences in candidates' beliefs about behavioral responses — rather than differences in their distributional preferences — may fully account for the observed differences in their tax proposals. The entire gap between Trump and Clinton's 2016 proposals, for instance, is consistent with a difference in ETI beliefs of just 8\% (0.25 versus 0.27), with no difference in underlying political weights required.

The rest of the paper is organized as follows. Section \ref{sec:approach} describes the inverse optimal tax approach. We next show how the model parameters, namely the distribution of income and the elasticity of taxable income, are calibrated in Section \ref{sec:calibration}. We summarize the tax policy proposals of U.S. presidential candidates in Section \ref{sec:candidate_policy} and use the inverse optimal tax approach to find the implied political weights in Section \ref{sec:estimated_weights}. Section \ref{sec:robustness} then offers an analysis of how the results are affected by alternative assumptions, while Section \ref{sec:invert_inverse} highlights how candidates may have different political weights or assumptions about taxpayer behavior and the implications of the latter. Finally, Section \ref{sec:conclude} concludes.

\section{The Inverse Optimal Tax Approach}
\label{sec:approach}

\subsection{Outline and Related Studies}

The optimal taxation model developed by \citet{DM1971}, and later refined by \citet{Saez2001}, formulates a process for solving for optimal income tax rates taking social welfare functions, elasticities with respect to earnings, and the empirical income distribution as given. Under certain assumptions, one can find the optimal tax rates for a variety of different social welfare functions, such as Rawlsian, Utilitarian, Progressive, or Conservative. Later, this approach was generalized to abstract social welfare weights by \citet{SS2016}.

The study of optimal taxation takes the social welfare function as given.  But this begs an important question: What are the social welfare function of those involved in making tax policy? Economists have attempted to answer this question in different contexts, using the inverse optimal taxation approach first developed by \citet{BS2012}. This approach, as the name suggests, uses the revealed-preferences and inverts the optimal taxation model, solving for social welfare weights while taking the elasticities of earnings, the income distribution, and tax policies as given. This methodology finds the social welfare weights that rationalize a given tax policy as optimal.  \citet{TH2017} derive a method for implementing inverse optimal tax in closed form.

\citet{BS2012} apply this method to estimate the welfare functions of the French tax-benefit system.  They find non-Paretian outcomes, in the form of negative welfare weights on some groups.  \citet{LW2016} use the methodology on U.S. tax policy from 1979 to 2010 and discuss the benefits and limitations of the inverse optimal-tax methodology in detail. Their results suggest a dilemma, threatening either the validity of the inverse optimum approach or the soundness of assumptions required for welfare analysis. \citet{Berg2019} implements the method using Norwegian data and highlights the concerns for horizontal and vertical equity that are implied by the estimated welfare weights.  Most recently, \citet{Embree2023} uses the inverse optimal tax method to find welfare weights of the tax systems across US states.

A closely related to paper to this study is \cite{JJZ2017}, who implement inverse optimal-tax methods to uncover political weights on parties in the Netherlands. The Netherlands is unique in that the Netherlands Bureau for Economic Policy Analysis conducts a detailed independent economic analysis of electoral policy proposals. The authors use this information to calculate welfare weights held by political parties for all income groups and unemployed. They verify some expected outcomes such as that left wing parties place a higher weight on the poor than right wing parties, and all political parties place a low weight on the rich. They also discuss the political nature of social welfare preferences, in that political biases and outcomes affect the nature of the weights. They derive the simplified formula used to calculate Social welfare weights.

Other related studies include the pioneering work has been done by \cite{BS2012}, who recover the social welfare weights for the tax system of France. Social welfare weights have also been recovered by  \cite{Petersen2007} for Denmark, \cite{BBHS2009} for single mothers in Germany and the UK, \cite{BK2010} for childless singles in Ireland and the UK over several decades, \cite{BDNPS2014} for childless singles in 17 European countries and the US, \cite{Hendren2014} for the US, \cite{LW2016} for the US over the period 1979–2010, and \cite{Lockwood2020} for the US if individuals have present bias. Finally, \cite{LS2020} use the inverse optimal-tax method to identify Pareto-improving tax reforms in Germany.
We extend this prior work by analyzing the political weights of candidates for President of the United States.

\subsection{Model}

\cite{AS2015} show that the optimal tax schedule can be characterized as a function of observed earnings, $z$ with the following condition:

\begin{equation}
\label{eqn:g_z}
  \frac{T'(z)}{1 - T'(z)} = \frac{1-F(z)}{\varepsilon z f(z)}\int_{z}^{\infty}\frac{1-g(\tilde{z})}{1-F(\tilde{z})}dF(\tilde{z})
\end{equation}

The optimality condition suggests that marginal tax rates, $T'(z)$ are inversely related to behavioral responses that induce deadweight loss (where $\varepsilon$ summarizes behavior through the compensated elasticity of taxable income) and inversely related to the social welfare weights placed on taxpayers with income $z$, $g(z)$.  The distribution of income, with CDF $F(z)$ is relevant to this calculation since increasing marginal tax rates on those with income $z$ implies higher taxes on all of those with income above $z$ as well.  In short, the optimal income tax schedule is one that raises a given amount of revenue while balancing the deadweight loss of taxation (efficiency concerns) with distributional preferences related through the social welfare function summarized by the weights, $g(z)$ (equity concerns).

Following \citet{JJZ2017}, with a constant elasticity of taxable income (ETI), and no income or participation effects, we can invert \ref{eqn:g_z} and solve for the marginal social welfare weights (MSWW), $g(z)$ as a function of marginal tax rates, the distribution of income, and the ETI.  Assuming $T'(z)$ is optimal, then the MSWWs can be written as:
\begin{equation}
\label{eqn:g_z}
    g_z = 1 + \theta_z \varepsilon^c \frac{T'(z)}{(1-T'(z))} + \varepsilon^c \frac{zT''(z)}{(1-T'(z))^2},
\end{equation}

where $\theta_z \equiv 1 + \frac{zf'(z)}{f(z)}$.  This equation says that the MSWW is equal to one plus the marginal deadweight loss of taxation.  That is, if we observe a planner imposing high marginal tax rates on taxpayers with income $z$, which result in a large increase in deadweight loss, then it must be that the planner's social welfare weight on those taxpayers is high. If the welfare weight were not high, then the planner could reduce the marginal tax rate, increasing efficiency with little cost to the social welfare. The marginal weight, $g(z)$, can be interpreted as the value of transferring one dollar of income to someone with income $z$.

\citet{JJZ2017} derive the expression for $g(z)$ in a more general model with income and participation effects.  However, they find that income and participation effects have quantitatively negligible impacts on the estimated MSWWs. We therefore use the more parsimonious model, which still captures the main determinants of the MSWWs.  Excluding income effects from our model will tend to bias downward the social welfare weights on high income earners since the presence of income effects would lower the distortion of high tax rates on labor supply, suggesting that optimal tax rates would be higher when there are income effects.  But given the evidence of \citet{JJZ2017} we do not expect the bias to be significant.  

There is evidence which suggests that the $\varepsilon^c$ varies across $z$ (see \citet{SSG2012}), with higher income taxpayers generally having a larger ETI. We relax this assumption in our analysis in Section \ref{sec:nonconstant_eti}, but here we note that the constant ETI assumption will tend to bias upward the MSWWs on high income earners.  This is because using a lower ETI for high income earners will reduce the deadweight loss implied from taxing them and therefore and therefore must mean that they are receiving a relatively higher MSWW given the relatively high marginal tax rates they face.

We follow the existing inverse optimal taxation literature and focus only on labor income.  Excluding capital income has an ambiguous effect on the MSWWs.  On the one hand, the distribution of capital income is more skewed towards the top earners, which means more revenue could be raised (and redistributed) if tax rates on these earners were higher.  On the other hand, the ETI for capital income is probably larger than that of labor income, which would tend to suggest lower tax rates on capital income. The implication for the optimal tax schedule, and therefore the implied social welfare weights from the inverse approach, is an empirical question about which effect dominates.

Finally, we refer to the MSWWs in our context as political weights.  Since we are analyzing the policy proposals of political candidates, it's not the case that we are estimating the candidates' social welfare functions directly.  Rather, we're estimating the weights they place on individuals as part of a policy position that they are relating in order to gain favor electorally. Hence, the weights implied by the inverse optimal tax method applied to the candidates' policies are more precisely defined as political weights rather than social welfare weights.

% There are a few special cases worth considering. First, the case in which $\varepsilon^c = 0 \implies g_z \equiv 1$, which gives us the usual utilitarian social welfare function. If $T'(z) \equiv 0$ in the case of lump sum taxes, this also implies that $T''(z) \equiv 0$, and thus $g_z = 1$, the utilitarian social welfare function. If marginal tax rates are constant (i.e., a flat tax), then $T''(z) \equiv 0$ and the
% implied weights come down to a function of how much revenue can be generated
% from different parts of the income distribution ($\theta_z$).

\section{Calibration}
\label{sec:calibration}

From Equation \ref{eqn:g_z}, it's clear that to recover the MSWWs, one needs to observe the income distribution, $f'(z)$, the marginal tax rate schedule, $T'(z)$, and have an estimate of the ETI, $\varepsilon$. In this section, we describe how we estimate or calibrate each of these objects.

\subsection{The Distribution of Earnings}

The distribution of earnings, $f(z)$, is a density function which represents the likelihood of a taxpayer having income, $z$. We approximate the earnings distribution with a Pareto log-normal distribution \citep{HajargashtGriffiths2013} and estimate the parameters of the distribution via maximum likelihood using a data file produced by the \texttt{taxdata} project. This file uses the 2012 IRS Public Use file as it's base.  It then imputes non-filers using the March Current Population Survey from 2012, resulting in about 200,000 observations.  These data are extrapolated out through 2024 (the last election year we consider) by growing out nominal variables at observed rates of growth and using a linear optimization routine that adjusts the sampling weights to high IRS and other targets from each year.  The results is a data file representative of tax units in the U.S. for all years, 2012-2024.

Figure \ref{fig:income_dist} shows the fitted Pareto log-normal distribution for 2012.  Each year we consider (2013-2024) looks roughly similar, with a large mass concentrated in lower income bins and a long right tail.

\begin{figure}[H]
\caption{Pareto Log-Normal Density Estimate of $f(z)$}
\label{fig:income_dist}
\centering
\includegraphics[width=\textwidth]{./images/income_dist_histogram.png}
\end{figure}

%We require the density function to be differentiable over the interval of estimation. Thus, we use Kernel Density Estimation (KDE), 
%\begin{equation*}
%    \hat{f}(z) = \frac{1}{h*n} \sum_{k=1}^{n} K(\frac{z-z_i}{h}),
% \end{equation*}

We differentiate $f(z)$ to find $f'(z)$.

\subsection{Marginal Tax Rates}\label{sec:mtr_est}

Using the open-source \texttt{Tax-Calculator} microsimulation model, we simulate each candidate's tax policy proposals and find the resulting marginal tax rates faced by the primary filers in each tax unit in our underlying data file.  We use the same tax return microdata described in the previous section for this analysis.  When computing marginal tax rates, we focus on the rate associated with earning another dollar of labor income (wages and salaries on IRS Form 1040) for the primary filer.  We also use the combined, income and payroll tax, marginal tax rate to reflect the full burden of the federal tax system on labor income.

After simulating the marginal tax rates for each tax unit in our sample under a particular policy, we find the distribution of marginal tax rates over income using a kernel regression.  Specifically, we collapse the underlying data into 1,000 bins based on labor income with equal dollar width.  For each bin, we compute the weighted average marginal tax rate, using the sampling weights from our microdata.  A kernel regression is then fit to these binned data. This gives a smooth function for $T'(z)$.  An example of the resulting marginal tax rate function, using 2023 tax law, can be seen in Figure \ref{fig:mtr_2023}. 

\begin{figure}[H]
\caption{Fitted Marginal Tax Rate Schedule, $T'(z)$, for 2023 law}
\label{fig:mtr_2023}
\centering
\includegraphics[width=\textwidth]{./images/MTR_2023.png}
\end{figure}

$T''(z)$ is found through numerical differentiation.

In Section \ref{sec:candidate_policy}, we describe the policies of the candidates we consider.  These policies are then parameterized and fed into the \texttt{Tax-Calculator} model to compute effective marginal tax rates, as described above.

% The primary model performs a kernel regression, also known as the Nadaraya–Watson estimator, which uses a multivariate kernel density estimation (univariate case described in section 4.1) of the joint density between $z$ and $T'$ to provide a local mean of $T'$ about $z$. We use a kernel bandwidth of \$10,000. The estimate is calculated as follows.

% \begin{equation*}
%     \hat{T'}(z) = \sum_{i=1}^{n} \frac{K(z-z_i)}{\sum_{i=1}^{n} K(z-z_i)} T'_i(z_i)
% \end{equation*}

% Where $K$ is the Gaussian Kernel (described in section 4.1), $n$ is the sample size, and $T'_i(z_i)$ is the marginal tax rate for individual $i$.
% Figure 5 is the kernel regression estimate of $T'(z)$ for each policy proposal.

% This kernel estimate is sufficiently smooth (continuous derivative), and accounts for the variation in policies. However, we recognize that the differences in the estimate vary much more across sample years than it does across policies. We discuss the implications of this result in sections 5 and 6. 

\subsection{Elasticity of Taxable Income}\label{sec:ETI}

\citet{SSG2012} survey the literature on the elasticity of taxable income $\varepsilon^c$. This is an important component of the model because it summarizes the behavioral responses to taxation. Behavioral responses to taxation lead to deadweight losses because the change in behavior results in a larger loss in private transactions (e.g., gains to employers and employees from work) than there is a gain in tax revenue.  Without behavioral responses, the revenue collected would exactly equal the losses to private parties and there would be no overall economic loss. 

The elasticity of taxable income operates as a sufficient statistic, summarizing a broad range of responses to taxes, such as the labor supply elasticity, timing of realization of capital gains, increased charitable expenses, and others. Seminal work by \citet{Feldstein1999} shows that the deadweight loss of taxation can be computed solely on the basis of the ETI. The elasticity depends on a number of factors including personal preferences and the tax system. A tax system which significantly encourages changing behavior, such as through many itemized deductions, will generate a higher $\varepsilon^c$ and therefore more deadweight loss. 

\cite{SSG2012} find that most estimates of average long run elasticity of taxable income range between 0.12 and 0.4, with a middle of the road estimate being 0.25. We use $\varepsilon=0.25$ in our baseline results, but show robustness to alternative assumptions in Section \ref{sec:robustness}.

\section{Tax Policies of Presidential Candidates}\label{sec:candidate_policy}
We estimate the political weights associated with the tax policy proposals of the Democrat and Republican general election candidates from 2012-2024.  In most cases, we use the policy proposals set forth in official campaign materials.   In a couple cases (Obama 2012 and Trump 2020) there was not a clear tax policy position from campaign materials. In those cases we've identified their policy positions public positions they've supported near those elections.

Below, we summarize the positions of each candidate that we model to compute the marginal tax rate schedule that results from their proposals.\footnote{The specific parameterization of these policies are available in the Github repository...} Note that we exclude business tax proposals and proposals that we lack data to estimate (e.g., elimination of 529 Plans proposed by Obama).  Thus, the proposed changes outlined below are the set of tax law changes proposed by the candidates that affect the marginal tax rates in the theoretical model (i.e., those on labor income) and which we can capture with our data and microsimulation model.

\subsubsection{2012 Election}

The 2012 Presidential Election major party candidates were Mitt Romney (Republican) and Barack Obama (Democrat). Relative to 2012 law, Romney proposed the following changes in tax policy:\footnote{See \citep{Romney} for details.}:
% https://taxfoundation.org/simulating-economic-effects-romneys-tax-plan/
\begin{itemize}
    \item {A 20 percent reduction in all personal marginal income tax rates}
    \item {Elimination of the capital gains Tax for low and middle income earners}
    \item {Eliminate the Alternative Minimum Tax}
    % \item Reduce Corporate Income Tax from 35 percent to 25 percent
    % \item Eliminate the federal Estate Tax
\end{itemize}

President Obama did not lay out a clear set of tax policy changes during the 2012 election, so we take his positions from the 2015 State of the Union Address.  We assume that the 2015 speech represents President Obama's views on the optimal tax system.  The changes proposed in this speech included :\footnote{See \citep{Obama} for details.}:

% https://taxfoundation.org/basics-president-obamas-state-union-tax-plan/#:~:text=According%20to%20reports%2C%20the%20president's,tax%20increases%20over%20ten%20years.&text=The%20largest%20tax%20increase%20in,tax%20rate%20to%2028%20percent.&text=Recent%20reports%20point%20towards%20a,24.2%20percent%20plus%203.8%20percent).

\begin{itemize}
    \item {\$500 tax credit for two earner families}
    \item {Increase the Child Tax Credit from \$1,050 to \$3,000, and increase the maximum income from \$43,000 to \$120,000}
    \item {Increase the EITC from \$503 to \$1,000 for taxpayers without children}
    % \item Exempts Pay-As-You-Earn Student Loan Forgiveness from income taxation
    \item Increase Capital Gains Tax Rate to 28 Percent
    % \item End ``Stepped-Up Basis''
    % \item Bank tax: .07 percent tax on liabilities of firms with \$50 billion in assets
    % \item Cap of \$3.4 million for retirement plans and IRAs
    % \item Eliminate 529 plans - tax exempt savings for education
    % \item 28 percent corporate income tax with minimum tax on foreign earned profits
\end{itemize}


\subsubsection{2016 Election}

The two major party candidates in the 2016 Presidential Election were Donald Trump (Republican) and Hillary Clinton (Democrat).  Both had significant tax policy proposals to change tax law relative to current law in 2016.  President Donald Trump proposed the following changes to United States tax policy while campaigning in 2016 \citep{Trump}:
\begin{itemize}
    \item Adjust the individual income tax schedule:
        \begin{center}
        \begin{tabular}{||c c c c||} 
         \hline
         Single Filer & Married Joint Filers &  Ordinary Income Rate & Capital Gains Rate\\ [0.5ex] 
         \hline\hline
         [\$0, \$37500) & [\$0, \$75000) & 12\% & 0\%\\ 
         \hline
         [\$37500, \$112500) & [\$75000, \$225000) & 25\% & 15\%\\ 
         \hline
         >\$112500 & >\$2250000 & 33\% & 20\%\\ 
         \hline
        \end{tabular}
        \end{center}
    \item Eliminate head-of-household filing status
    \item Eliminate the Net Investment Income Tax
    \item Increase the standard deduction from \$6,300 to \$15,000 (singles)
    \item Make childcare costs deductible for low and middle income earners, up to the average cost
    \item Create credits for childcare of \$1,200 through the EITC
    \item Cap itemized deductions at \$100,000 (singles)
    \item Eliminate individual and corporate alternative minimum tax
    % \item Tax carried interest as ordinary income
    % \item Reduce corporate income tax from 35\% to 15\%
    % \item Allow manufacturing firms to choose between full expensing of capital investment and the deductibility of interest paid
    % \item Eliminate all business credits except research and development credit
    % \item Expand credits for employer provided day care
\end{itemize}
%see: https://taxfoundation.org/details-analysis-donald-trump-tax-plan-2016/

%and the JSON we are using (to see what is modeled here): https://github.com/PSLmodels/examples/blob/main/psl_examples/taxcalc/Trump2016.json

Former Secretary of State Hillary Clinton proposed the following changes during her presidential campaign in 2016 as the Democratic nominee (changes relative to 2016 law)\footnote{See \citep{Clinton} for details.}:
% https://taxfoundation.org/details-and-analysis-hillary-clinton-s-tax-proposals-october-2016/
\begin{itemize}
    \item Create 4\% ``surcharge'' on AGI above \$5 million
    \item Create a 30\% minimum tax rate on AGI above \$1 million, phased in up to \$2 million
    \item Limit the value of itemized deductions (except charity) and tax expenditures to a tax rate of 28\%
    \item Increase Child Tax Credit by \$1,000 for children under 5, phase in from first dollar rather than \$3,000
    % \item Create new schedule of long term capital gains tax rates:
    
    % \begin{center}
    % \begin{tabular}{||c c||} 
    %  \hline
    %  Years Held & Marginal Tax Rate \\ [0.5ex] 
    %  \hline\hline
    %  [0, 2) & 39.6\% \\ 
    %  \hline
    %  [2, 3) & 36\% \\
    %  \hline
    %  [3, 4) & 32\% \\
    %  \hline
    %  [4, 5) & 28\\
    %  \hline
    %  [5, 6) & 24\\  
    %  \hline
    %  6+     &  20\\
    %  \hline
    % \end{tabular}
    % \end{center}
    
    % \item Limit value of tax-deferred and tax-free savings (how much?)
    % \item Limit tax benefit of like-kind exchanges
    % \item Tax carried interest at income tax rate rather than capital gains rate
    % \item 20\% credit for caregiver expenses and \$5,000 in tax relief for excessive healthcare costs
    % \item Establish 'Risk Fee' on large banks
    % \item Small business benefits:
    % \begin{enumerate}
    %     \item Expand expensing
    %     \item Increase business eligible for cash accounting
    %     \item Increase start up deduction from \$5,000 to \$20,000
    %     \item Create small business standard deduction
    %     \item Expand ACA credit for small businesses
    % \end{enumerate}
    % \item Eliminate fossil fuel tax expenditures
    % \item Establish business tax credit for profit sharing and apprenticeships
    % \item Decrease estate tax exemption to \$3.5 million, introduce progressive rates with top rate of 65\% for estates worth \$500 million or more
    % \item Tax unrealized gains on death, with some exemptions
    % \item Tax on high frequency trading
\end{itemize}


\subsubsection{2020 Election}
The election of 2020 featured President Donal Trump (Republican) and President Joe Biden (Democrat).  The Trump campaign did not articulate a set of tax policy proposals, so we take his signature piece of legislation during his first term, the Tax Cuts and Jobs Act, as his tax policy position while campaigning for re-election.  Relative to 2017 law, the TCJA made the following changes relevant to the MTRs we model: \footnote{See \citep{TCJA} for details.}
\begin{itemize}
    \item Alter the individual marginal income tax schedule:
    \begin{center}
    \begin{tabular}{||c c||} 
     \hline
     Income  & Marginal Tax Rate \\ [0.5ex] 
     \hline\hline
     [0, \$9,525) & 10\% \\ 
     \hline
     [\$9,525, \$38,700) & 12\% \\
     \hline
     [\$38,700, \$82,500) & 22\% \\
     \hline
     [\$82,500, 157,500) & 24\% \\
     \hline
     [\$157,000, \$200,000) & 32\% \\  
     \hline
     [\$200,000, 500,000) & 35\% \\
     \hline
     >\$500,000 & 37\% \\
     \hline
    \end{tabular}
    \end{center}
    \item Index these changes using Chained CPI
    % \item The jointly filing schedule has double the income intervals except with the top rate starting at \$600,000
    \item Increase the standard deduction to \$12,000 for single filers, \$18,000 for heads of household, and \$24,000 for joint filers
    \item Eliminate the personal exemption 
    \item Double the Child Tax Credit and increase the phaseout from \$110,000 to \$400,000
    \item Limit a number of deductions such as the mortgage interest deduction and the state and local tax deduction
    \item Raise the exemption on the alternative minimum tax from \$109,400 for married filers, increase the phaseout to \$1 million
    % \item Double the estate tax exemption
    % \item Lower the corporate tax rate to 21\%
    % \item Eliminate the corporate alternative minimum tax
    % \item Limit several credits and the deductibility of various business expenses such as net interest expense, domestic production activities deduction, orphan drug credit and rehabilitation credit
\end{itemize}


President Biden proposed the following tax policy changes on labor income (relative to 2020 law)\footnote{See \citep{Biden} for details.}:
%https://taxfoundation.org/joe-biden-tax-plan-2020/#Details
\begin{itemize}
    \item Impose a 12.4\% Social Security Payroll tax on income beyond \$400,000
    \item Increase the top marginal income tax rate from 37\% to 39.6\%
    \item Tax Capital Gains at the income tax rate of 39.6\% on income beyond \$1 million
    \item Cap the benefit of itemized deductions to 28 percent of value and restores the ``Pease Limitation'' for those with income greater than \$400,000
    \item Phase out the qualified business income deduction for filers with taxable income above \$400,000
    \item Expand EITC for childless workers 65 and older
    \item Expand the Child and Dependent Care Tax Credit from a maximum of \$3,000 to \$8,000 and increase the maximum reimbursement rate from 35 percent to 50 percent
    % \item Re-establish First-Time Homebuyers’ Tax Credit of \$15,000
    % \item Restores the Estate Tax to 2009 levels
    % \item Increases the corporate income tax rate from 21 percent to 28 percent
    % \item Creates a minimum tax on corporations with profits of \$100 million or higher
    % \item Double the tax rate on Global Intangible Low Tax Income earned by foreign subsidiaries of US firms from 10.5 percent to 21 percent
    % \item Establish a Manufacturing Communities Tax Credit
    % \item Expand the New Markets Tax Credit and makes it permanent
    % \item Offer tax credits to small business for adopting workplace retirement savings plans
    % \item Expand several renewable-energy-related tax credits
\end{itemize}


\subsubsection{2024 Election}

The 2024 election featured President Donald Trump (Republican) and Vice President Kamala Harris (Democrat).  President Trump proposed the following changes to individual federal income taxes (relative to 2024 law):

\begin{itemize}
 \item Extend the individual provisions of the Tax Cuts and Jobs Act
 \item Eliminate taxes on Social Security Benefits
 % \item No taxation of tip income
 % \item Extend the business tax provisions of the TCJA
 % \item Lower the corporate income tax rates to 15 percent
\end{itemize}

Vice President Harris proposed the following changes to individual federal income taxes (relative to 2024 law):

\begin{itemize}
 \item Expand the child tax credit and make it fully refundable
 \item Expand the Earned Income Tax Credit
 % \item No taxation of tip income
 % \item Provide a \$6,000 refundable tax credit for newborns.
 % \item Permanently extend enhanced premium tax credits
 % \item Raise the corporate income tax rate to 28\%
\end{itemize}

Both candidates also proposed exempting tip income from taxation, but we do not have data on tip income to model this proposal and so we exclude it and Harris' proposal for a refundable credit for newborns.



\section{Political Weights of Recent Presidential Candidates}
\label{sec:estimated_weights}

Given the candidate proposals outlined in the previous section, and using the kernel regression estimator described in Section \ref{sec:mtr_est}, we estimate the marginal tax rate schedule for each of the eight candidate-platforms.  Figure \ref{fig:mtr_all} plots these, with blue lines denoting Democrat party candidates and red lines denoting Republican party candidates. Line type varies with the year of the election. 

\begin{figure}[H]
\caption{Candidates' Proposed Marginal Tax Rates, $T'(z)$}
\label{fig:mtr_all}
\centering
\includegraphics[width=\textwidth]{./images/mtr_all.png}
\end{figure}

A few observations are clear from Figure \ref{fig:mtr_all}. First, marginal tax rates for all candidates are very similar for incomes up to about \$100,000.  Second, Romney is an outlier with significantly lower marginal tax rates than other candidates, all along the income distribution.  Third, all candidates' proposals reflect a progressive marginal tax rate structure, although it is very modest with Romney.  Fourth, we see a clear tendency for Democratic candidates to have higher tax rates on the high income, than Republican candidates.  This would not be surprising to anyone who has observed US politics.  Finally, there does seem to be a trend over of time, and across parties, to more heavily tax the highest income earners.  Both Biden and Trump (2020) have the highest marginal tax rates of any general election candidate from their party since 2012.  Biden's plan to impose payroll taxes on incomes above \$400,000 shows up very clearly in these marginal tax rate functions, with a steep increase in the level of MTRs beginning around that threshold. This will be reflected in the implied political weights.

Combining the estimates of $T'(z)$ with our estimates on the income distribution in 2012, 2016, 2020, and 2024 using an ETI of 0.25, we back out the implied political weights with the inverse optimal tax approach, as given in Equation \ref{eqn:g_z}.  Figure \ref{fig:gz_all} relates the estimated weights.

\begin{figure}[H]
\caption{Estimated Political Weights, $g_z$}
\label{fig:gz_all}
\centering
\includegraphics[width=\textwidth]{./images/gz_all_pln.png}
\end{figure}

We note that for incomes below about \$200,000, there is little deviation among the candidates as to the weights they place on taxpayers.  Furthermore, weights are generally declining in income, across all candidates, showing a political motive to redistribute towards lower income taxpayers.   Candidates' political weights become more different at incomes over \$600,000.  At that point, we begin to see lower weights from Democratic candidates relative to Republicans.  At the same time, the pattern of weights within party are quite similar over time.  Also of note, weights are never negative.  A negative weight would indicate a non-Paretian policy by implying that marginal tax rates are set such that one is on the right side of the Laffer Curve, and a reduction in those rates would raise those agents utility and raise revenue (that could be redistributed). 

TODO: more to write here?  Be nice if we can get a plot to decompose the weights to work and put that here (we only have decomposition code for JJZ method but we are using the LW method to estimate the weights.  E.g.,

\begin{figure}[H]
\caption{Decomposing Political Weights, $g_z$, Biden 2020}
\label{fig:Biden_decomp}
\centering
\includegraphics[width=\textwidth]{./images/composition_Biden2020_gz.png}
\end{figure} 


\section{Robustness}
\label{sec:robustness}

\subsection{Sensitivity to the level of the ETI}
Figure \ref{fig:gz_vary_eti} shows the political weights from the Biden 2020 proposals under alternative assumptions about the (constant) ETI.  One can see that as the ETI increases, the implied political weights go down and that this pattern is more significant at the top of the income distribution where marginal tax rates are higher.


\begin{figure}[H]
\caption{Political Weights, $g_z$, for Biden 2020, alternative $\varepsilon$}
\label{fig:gz_Biden_vary_eti}
\centering
\includegraphics[width=\textwidth]{./images/vary_ETI_Biden2020_gz_numerical.png}
\end{figure} 

\subsection{Relaxing the assumption of a constant ETI}\label{sec:nonconstant_eti}

Here, we relax the assumption that the ETI be constant across the income distribution. Note that the condition in Equation \ref{eqn:g_z} relied on the assumption of constant ETI for its derivation. Following \citet{LW2016}, the first order condition for the planner is

\begin{equation}
\label{eqn:FOC}
    g_z f(z) + \frac{d}{dz}\left [
    (1-F(z)) - \varepsilon(z) \frac{T'(z)}{1-T'(z)} zf(z)
    \right ] = 0
\end{equation}

Relaxing the constant ETI assumption and applying the chain rule, we obtain Equation \ref{eqn:g_z} with an additional term depending on $\varepsilon'(z)$:

\begin{equation}
\label{eqn:g_z_vary_ETI}
    g_z = 1 +  \varepsilon(z) \left [ \theta_z \frac{T'(z)}{(1-T'(z))} + \frac{zT''(z)}{(1-T'(z))^2} \right ] + \varepsilon'(z) \frac{z T'(z)}{1- T'(z)}.
\end{equation}

To calibrate the value at different points in the income distribution, we turn to \citet{GS2002}, who report the value of the ETI for three income groups.  We take their estimated ETI values from Table 9 and assign them to the income level that is approximately the median of each of the reported income bins (\$30,000, \$75,000, and \$250,000, respectively).  These are represented with red dots in Figure \ref{fig:vary_eti}.  Next, we take the ETI of the top 1\% estimated by \citet{MW2000} to be 1.83, using this as the value of the ETI at \$2,000,000 of income. We set another value of 1.9 at \$10,000,000 so that the spline function with estimates has a more flat curvature at high income levels.  The result of fitting a cubic spline through these data points is a function represented by the blue line in Figure \ref{fig:vary_eti}. 

\begin{figure}[H]
\caption{Varying $\varepsilon$, fitted values from Gruber and Saez (2002)}
\label{fig:vary_eti}
\centering
\includegraphics[width=\textwidth]{./images/ETI_spline.png}
\end{figure} 

With the function $\varepsilon(z)$ from Figure \ref{fig:vary_eti}, we again recover the political weights impied by each of the eight candidates' platforms.  The results are in Figure \ref{fig:gz_vary_eti}.  Here we see some negative weights and generally a much steeper decline in the implied political weights on high income earners.  Of course, we noted earlier that a higher ETI will result in lower weights as the welfare function much be placing less weight on transfers to taxpayers at a given point in the income distribution if the deadweight loss is very large yet these filers face high marginal tax rates.


\begin{figure}[H]
\caption{Political Weights, $g_z$, $\varepsilon(z)$}
\label{fig:gz_vary_eti}
\centering
\includegraphics[width=\textwidth]{./images/gz_numerical_all_vary_eti.png}
\end{figure} 

\section{Inverting the inverse}
\label{sec:invert_inverse}

Here we invert the invert the inverse.  That is, we use Equation \ref{eqn:FOC}, but solve for $\varepsilon(z)$ as function of the $T'(z)$, $f(z)$, and $g(z)$. Given that there is considerable debate among economists about the elasticity of taxable income \citep{SSG2012}, can we expect policy makers to agree as well? 

The thought experiment we want to do with this is the following: we see two candidates different policies.  If the median voter theorem holds, they should converge on their policy preferences.  So perhaps it's not a disagreement about the political preferences, but rather differences over beliefs about behavioral responses to policy as captured by the Elasticity of Taxable Income in the model.

This is best solved as recasting Equation \ref{eqn:g_z_vary_ETI} as a first-order linear ODE in terms of $\varepsilon(z)$. The solution thus requires a boundary condition, which we set $\varepsilon_0$ equal to the estimate in Section \ref{sec:ETI}. With Equation \ref{eqn:solve_eti}, we can say solve for the implied beliefs about the ETI given a set of political weights and a proposed tax policy. 

\begin{equation}
    \label{eqn:solve_eti}
    \varepsilon(z) = \frac{1-T'(z)}{T'(z) z f(z)}\left[ 
    \frac{\varepsilon_0 T'(z_0) z_0 f(z_0)}{1-T'(z_0)} - \int_0^z (1-g_z(\tilde z)) f(\tilde z)d \tilde z
    \right]
\end{equation}

%\begin{equation}\label{eqn:solve_eti2}
%\varepsilon_{z} = \frac{(1-T'(z))}{T'(z)}\frac{(1-F(z))}{zf(z)}\int_{z}^{\infty}\frac{1-g_(\tilde{z})}{1-F(z)}dF(\tilde{z})
%\end{equation}

By assuming a set of candidates has the same political weights across the distribution of income, we can measure what their beliefs about the ETI must be to justify their policies. We illustrate this with Trump and Clinton 2016.  The implied political weights from their policy proposal and $\varepsilon=0.25$ are shown in Figure \ref{fig:gz_TrumpClinton}.

\begin{figure}[H]
\caption{Political Weights, $g_z$, Clinton and Trump 2016}
\label{fig:gz_TrumpClinton}
\centering
\includegraphics[width=\textwidth]{./images/trump_clinton_g_z_numerical.png}
\end{figure} 

So if Trump really had the same political preferences at Clinton (and she believes $\varepsilon=0.25$), then what must his beliefs about the ETI be?  Figure \ref{fig:trump_eti} shows this. We see that we can rationalize the differences in policy between Trump and Clinton's policies by attributing a belief to Trump of $\varepsilon_z \approx 0.27$ rather than $\varepsilon_z \approx 0.25$. This is well within the reasonable range of $(0.12, 0.4)$ discussed by \citet{SSG2012}, and thus could be a reasonable explanation of the policy difference. TODO: Is \citet{Hendren2020} at good cite here?


\begin{figure}[H]
\caption{Implied $\varepsilon(z)$ for Trump 2016}
\label{fig:trump_eti}
\centering
\includegraphics[width=\textwidth]{images/eti_trump.png}
\end{figure} 


We can also use this method to rule out preferences which would be associated with wildly counterfactual beliefs about the ETI for a given set of tax proposals. For example, if we assume all weights are utilitarian (i.e. $g(z) \equiv 1$), what beliefs would candidates need to have to justify their proposals? Figure \ref{fig:eti_utilitarian} shows this. Candidates would need to believe an ETI ranging from $\varepsilon_z \in (0.07, 0.11)$ for most of the income distribution, which lies outside the credible range of estimates. Thus, we can reject the hypothesis that presidential candidates in the US are plausibly utilitarian. 

\begin{figure}[H]
\caption{Implied $\varepsilon(z)$ for all candidates assuming utilitarian weights}
\label{fig:eti_utilitarian}
\centering
\includegraphics[width=\textwidth]{images/eti_utilitarian.png}
\end{figure} 


%\begin{figure}[H]
%\caption{Implied %$\varepsilon(z)$ for Trump 2016}
%\label{fig:trump_eti}
%centering
%includegraphics[width=\textwidth]{./images/trump_eti.png}
%\end{figure} 


\section{Conclusions}
\label{sec:conclude}

\subsection{Summary}
We use the derivations of \citet{JJZ2017} to implement the inverse-optimal tax methodology, to infer the social welfare weights placed by presidential candidates on individuals in society. We use the IRS Public Use files and the \texttt{Tax Calculator} microsimulation model to estimate the distribution of income and marginal tax rates. 

We use these inputs to calculate the political weights which make the proposed tax policies optimal for each presidential candidate. We find that candidates from both parties place similar weights on taxpayers with incomes below \$100,000, but that the weights of Democrat candidates on taxpayers with earnings above \$400,000 are much lower than Republicans' weights on these taxpayers.  We also see some evidence that both parties place lower weight on these high income taxpayers in 2020 than in prior election cycles.

Under the assumption of a constant ETI of 0.25, all estimated weights are positive, indicating Paretian policies from the eight candidates we analyze.  However, when we use an empirically founded ETI function that varying by income, we do find negative welfare weights on high income taxpayers. This suggests that candidates may not believe that high income taxpayers have significantly larger behavioral responses than lower income taxpayers.

Finally, we show results from an exercise where we try to see what ETI assumptions it would take for the Trump 2016 tax policy to reflect the same political weights as the Clinton 2016 tax policy (under the assumption that Clinton believes the ETI is constant across income at 0.25).  This exercise yields implied ETIs that are unreasonable (e.g., the wrong sign).  This provides some evidence that the underlying political weights of the parties differ, which runs against the idea of platform convergence implied by the median voter theorem.


\subsection{Future Work}
There are many avenues for future development and improvements in this project. The first of which is to better capture the right tail of the income distribution. A potential solution is to fit a Pareto distribution beyond a certain threshold following \citet{Saez2001}. We can accomplish this by using the parametric log-normal estimate and its derivative upt to a certain point in the income distribution and then solve analytically for parameters of the Pareto distribution above that point. The Pareto distribution is often used for modeling the upper tail of an income distribution because it has a meaningful interpretation in that the average beyond a threshold is proportional to the threshold. Figure \ref{fig:upper_tail} shows that for the 2021 data, the average of the upper tail over the threshold converges asymptotically, which implies that the upper tail is well represented by a Pareto distribution. Since the Pareto density is a strictly decreasing differentiable function, this would eliminate the oscillations and provide an analytical solution to the value of the density derivative where it is implemented \citep{CharpentierFlachaire2019}.

\begin{figure}[H]
\caption{Earnings Distribution Upper Tail Ratio, 2021}
\label{fig:upper_tail}
\centering
\includegraphics[width=\textwidth]{./images/Saezfig2_2021.png}
\end{figure}

Another addition we could make is to incorporate taxes of other sources of income, such as capital income, and wealth taxes. There are important differences across parties in candidates' policy proposals related to capital income taxation.  Wealth taxes are a hot topic of debate both academically and politically, with high profile debates being spurred by politicians such as Bernie Sanders and Elizabeth Warren. It would be interesting to be able to fulling examine those candidates' proposals with our methods.  \citet{SLRSZ2022} may provide a framework to use when considering the taxation of multiple income sources.

We could also improve the information included in our estimates by incorporating other transfer programs outside the tax system such as the Supplemental Nutrition Assistance Program (SNAP). The differences in policy goals between candidates is often divided along party lines. This could result in the candidates being more significantly different from each other. In particular, it might make the weights on the lower end of the income distribution higher since most transfer programs are need-based, and it would likely be higher for democrats than Republicans.  The \texttt{PolicyEngine-US} model provides the ability to compute effective marginal tax rates on US households that consider both tax and benefit programs.


Finally, we could expand our analysis to more candidates, such as those in third parties and those who did not win their party's nomination. We could also try to extend our analysis further into the past. For transparency and reproducibility, the model used for our analysis is open source and available on Github (at \href{https://github.com/PSLmodels/InverseOptimalTax}{PSLmodels/InverseOptimalTax}). Thus, all that is needed to extend our results to a new candidate or policy is the specific changes suggested in the proposal. This tool can be used for candidates in future elections as well.


\newpage


% -----------------------------
\bibliographystyle{apalike} % Set the bibliography style
\bibliography{references.bib} % Create the list of references

\newpage 

\section{Appendix}

\subsection{Optimal Tax Formula}
Define the ETI at income $z$ as
$$\varepsilon(z) = \frac{1-T'(z)}{z} \frac{\partial z}{\partial (1-T'(z))}$$
and consider an increase in the marginal tax rate $T'(z)$ by $d\tau$ on a narrow income band $[z, z+dz]$. Note that this increases the tax liability for everyone above $z+dz$ by $d\tau \cdot dz$. The fraction of the population to whom this applies is $1-F(z)$. So the mechanical revenue increase from this reform is
$$dR_1 = (1-F(z)) \cdot d\tau \cdot dz.$$
 
However, some people in the band $[z, z+dz]$ change their behavior due to the different marginal rate. There are $f(z)dz$ taxpayers in this band. From the definition of the ETI, the behavioral response in earnings for each of these taxpayers is
$$dz^* = -\varepsilon(z) \cdot \frac{z}{1-T'(z)} \cdot d\tau.$$
 
Each dollar of reduced earnings costs the government $T'(z)$ in lost revenue. Thus, the total behavioral revenue loss is
$$dR_2 = -\varepsilon(z) \cdot \frac{T'(z)}{1-T'(z)} \cdot z \cdot f(z) \cdot d\tau \cdot dz.$$
 
Finally, we must account for the welfare effect of the reform. The mechanical tax increase of $d\tau \cdot dz$ applies to all taxpayers above $z$. The social value of taking a dollar from a taxpayer at income $\tilde{z}$ is $g(\tilde{z})$, the marginal social welfare weight. Therefore, the welfare cost of the mechanical revenue increase is
$$dW = -\int_{z}^{\infty} g(\tilde{z}) \, dF(\tilde{z}) \cdot d\tau \cdot dz.$$

Note that we do not need to account separately for the welfare effect of the behavioral response, since the individual is optimizing and therefore a small change in $z$ has no first-order effect on utility by the envelope theorem. The behavioral response matters only through its fiscal externality, captured in $dR_2$.
 
At the optimum, the net welfare effect of the perturbation must be zero. The government values revenue at the margin at rate 1 (since it can redistribute it optimally), so the optimality condition is
$$dR_1 + dR_2 + dW = 0.$$
 
Substituting and dividing through by $d\tau \cdot dz$:
$$(1-F(z)) - \varepsilon(z) \frac{T'(z)}{1-T'(z)} z f(z) - \int_{z}^{\infty} g(\tilde{z}) \, dF(\tilde{z}) = 0.$$
 
We can rewrite the integral term as
$$\int_{z}^{\infty} g(\tilde{z}) \, dF(\tilde{z}) = (1-F(z)) - \int_{z}^{\infty} (1-g(\tilde{z})) \, dF(\tilde{z}).$$
 
Substituting and solving for the tax rate:
$$\varepsilon(z) \frac{T'(z)}{1-T'(z)} z f(z) = \int_{z}^{\infty} (1-g(\tilde{z})) \, dF(\tilde{z}),$$
 
which yields the optimal tax formula:
\begin{equation}
\label{eqn:optimal_tax}
\frac{T'(z)}{1-T'(z)} = \frac{1-F(z)}{\varepsilon(z) \, z \, f(z)} \int_{z}^{\infty} \frac{1-g(\tilde{z})}{1-F(z)} \, dF(\tilde{z}).
\end{equation}
 
This is the standard result from \citet{Saez2001} and Diamond (1998, NEEDS CITATION). When the ETI $\varepsilon(z) \equiv \varepsilon$ is constant across the income distribution, the formula simplifies to Equation \ref{eqn:g_z} in the main text.

\subsection{Inverse Optimal Tax Formula}

We now invert the optimal tax formula to recover the marginal social welfare weights implied by an observed tax schedule, following \citet{JJZ2017}. Taking the optimality condition from the previous section and writing it in integral form:

\begin{equation}
\int_{z}^{\infty} g(\tilde{z}) \, dF(\tilde{z}) = (1-F(z)) - \varepsilon(z) \frac{T'(z)}{1-T'(z)} z f(z).
\end{equation}

Define the right-hand side as
$$H(z) \equiv (1-F(z)) - \varepsilon(z) \frac{T'(z)}{1-T'(z)} z f(z).$$

Note that $H(z)$ captures the fiscal surplus from taxpayers above $z$, netting out the behavioral distortions. The first term is the mechanical tax base above $z$ (per unit of rate increase), and the second is the revenue lost through the behavioral response.

Differentiating the above wrt $z$:

$$-g(z) f(z) = H'(z) = -f(z) - \frac{d}{dz}\left[\varepsilon(z) \frac{T'(z)}{1-T'(z)} z f(z)\right].$$

Dividing by $-f(z)$ and rearranging yields the formula in \citet{LW2016}:
\begin{equation}
\label{eqn:gz_general}
g(z) = 1 + \frac{1}{f(z)}\frac{d}{dz}\left[\varepsilon(z) \frac{T'(z)}{1-T'(z)} z f(z)\right].
\end{equation}

When $\varepsilon(z)$ is constant in $z$, we evaluate the derivative to get the form in \citet{JJZ2017}:

\begin{align*}
g(z) &= 1 + \varepsilon \left[\frac{T''(z) z}{(1-T'(z))^2} + \frac{T'(z)}{1-T'(z)}\left(1 + \frac{z f'(z)}{f(z)}\right)\right].
\end{align*}

\subsection{Solving for $\varepsilon(z)$}

We start from the above Equation \ref{eqn:gz_general}. Since $\frac{d}{dz}(1-F(z)) = -f(z)$, this simplifies to:
\begin{equation*}
(g(z) - 1) f(z) = \frac{d}{dz}\left[\varepsilon(z) \frac{T'(z)}{1-T'(z)} z f(z)\right].
\end{equation*}

This is a first-order linear ODE in $\varepsilon(z)$. We solve it under two different boundary conditions.

\subsubsection{Solution with Initial Condition}
 
Integrate both sides of Equation of the above from $z_0$ to $z$:
$$\int_{z_0}^{z} (g(\tilde{z}) - 1) f(\tilde{z}) \, d\tilde{z} = \varepsilon(z) \frac{T'(z)}{1-T'(z)} z f(z) - \varepsilon(z_0) \frac{T'(z_0)}{1-T'(z_0)} z_0 f(z_0).$$

Solving for $\varepsilon(z)$:
\begin{equation*}
\varepsilon(z) = \frac{1-T'(z)}{T'(z) \, z \, f(z)}\left[\frac{\varepsilon_0 \, T'(z_0) \, z_0 \, f(z_0)}{1-T'(z_0)} - \int_{z_0}^{z} (1-g(\tilde{z})) f(\tilde{z}) \, d\tilde{z}\right]
\end{equation*}

\subsubsection{Solution with Transversality Condition}
 
Alternatively, we can integrate the above differential equation from $z$ to $\infty$. The transversality condition requires that the tax distortion term vanishes at infinity:
$$\lim_{z \to \infty} \varepsilon(z) \frac{T'(z)}{1-T'(z)} z f(z) = 0.$$

This is a standard assumption in the literature. Given a distribution with Pareto tail parameter $>1$ (which we estimate in the data), this holds as long as $\varepsilon(z)$ doesn't grow too fast and $T'(z)$ doesn't approach 0 or 1 too quickly in the limit. 

Integrating from $z$ to $\infty$:
$$\varepsilon(z) \frac{T'(z)}{1-T'(z)} z f(z) = \int_{z}^{\infty} (g(\tilde{z}) - 1) f(\tilde{z}) \, d\tilde{z},$$

Solving for $\varepsilon(z)$:
\begin{equation*}
\varepsilon(z) = \frac{1-T'(z)}{T'(z)} \cdot \frac{1}{z \, f(z)} \cdot \int_{z}^{\infty} (1-g(\tilde{z})) \, dF(\tilde{z}).
\end{equation*}

Note the parallel to the Diamond-Saez optimal tax formula. This version can also just be more simply solved for directly from the optimal tax formula by just solving for $\varepsilon(z)$.

\subsubsection{Comparing the Two Solutions}
When implementing these different approaches, we consider the simple thought experiment performed above to compare the two. A logical benchmark is that, holding weights $g(z)$ constant, uniformly lower marginal tax rates (as we see with Trump compared with Clinton in 2016) should imply a higher elasticity. This must be the case because the marginal value of public funds is assumed to be the same across the two scenarios (by the normalization of $g(z)$), so lower marginal tax rates can only be justified by more conservative beliefs about the distortions caused by those tax rates. However, the transversality formulation fails in this benchmark. In this example, for the initial condition, we see a uniformly higher implied ETI than that used to generate Clinton's $g(z)$. We see the opposite for the transversality condition: $\varepsilon(z) < 0.25$. This calls into question the veracity or aptness of this approach for the purpose of this experiment. Thus, we adopt the approach of the initial boundary condition at $\varepsilon_0$.

\begin{figure}[H]
\caption{Implied $\varepsilon(z)$ for Trump with Clinton's $g(z)$ under different boundary conditions}
\label{fig:eti_trump_compare}
\centering
\includegraphics[width=\textwidth]{images/eti_trump_compare_bc.png}
\end{figure} 


\end{document}